\documentclass[12pt, a4paper]{article}
\usepackage[T1]{fontenc}
\usepackage[utf8]{inputenc}
\usepackage[polish]{babel}
\usepackage[margin=2.5cm]{geometry}
\usepackage{array}
\usepackage{tabularx}
\usepackage{setspace}
\usepackage{verbatim}
\usepackage[hidelinks]{hyperref}
\usepackage{tocloft}
\usepackage{float}
\usepackage{listings}

\renewcommand{\cftsecleader}{\cftdotfill{\cftdotsep}}
\onehalfspacing

% Styl dla fragmentów kodu
\lstset{
    basicstyle=\ttfamily\small,
    breaklines=true,
    frame=single,
    
}

\title{\vspace*{3cm} \textbf{Dokumentacja Implementacyjna}\\Wizualizacja Grafów Planarnych (Java)}}
\author{Autorzy: Mateusz Bombola \\ Kajetan Karwacki}
\date{\today}

\begin{document}

\maketitle
\thispagestyle{empty}

\newpage
\tableofcontents
\newpage

\section{Architektura Aplikacji}
Program został zaimplementowany w języku Java z wykorzystaniem biblioteki \textbf{Swing}. Architektura opiera się na podziale na model danych, panel renderujący (View) oraz ramkę główną pełniącą rolę kontrolera (Controller).

\subsection{Główne klasy i struktury danych}
\begin{itemize}
    \item \texttt{Node}: Prosta klasa modelu przechowująca identyfikator (\texttt{int id}) oraz współrzędne (\texttt{double x, y}).
    \item \texttt{Edge}: Przechowuje nazwę krawędzi, identyfikatory wierzchołków incydentnych oraz wagę (\texttt{double weight}).
    \item \texttt{GraphPanel}: Klasa dziedzicząca po \texttt{JPanel}, odpowiedzialna za całą logikę graficzną i interakcje myszą.
    \item \texttt{GraphApp}: Główna klasa aplikacji (\texttt{JFrame}), zarządzająca menu, panelem bocznym i komunikacją z procesami zewnętrznymi.
\end{itemize}

\section{Szczegóły Implementacyjne Modułów}

\subsection{Zarządzanie stanem grafu}
Aplikacja wykorzystuje kolekcje z pakietu \texttt{java.util} do efektywnego zarządzania danymi:
\begin{itemize}
    \item \texttt{Map<Integer, Node> nodes}: Przechowywanie węzłów, gdzie kluczem jest ich ID. Pozwala to na szybki dostęp ($O(1)$) podczas rysowania krawędzi.
    \item \texttt{Map<Integer, Point2D.Double> snapshots}: Przechowuje pierwotne pozycje węzłów po wczytaniu współrzędnych. Funkcja \texttt{resetNodePositions()} wykorzystuje tę mapę do przywrócenia stanu grafu po ręcznych modyfikacjach użytkownika.
\end{itemize}

\subsection{Renderowanie i Transformacje}
W klasie \texttt{GraphPanel} zastosowano zaawansowane renderowanie przy użyciu \texttt{Graphics2D}:
\begin{itemize}
    \item \textbf{AffineTransform}: Wykorzystywany do operacji \texttt{translate} i \texttt{scale}. Pozwala to na separację współrzędnych świata (węzłów) od współrzędnych ekranu.
    \item \textbf{Antyaliasing}: Włączony poprzez \texttt{RenderingHints}, co zapewnia wysoką jakość rysowanych linii i kół.
    \item \textbf{Dynamiczne kolorowanie krawędzi}: Kolor krawędzi jest obliczany na podstawie jej wagi względem ekstremów (min/max weight) wczytanego grafu. Wykorzystywana jest składowa czerwona (R) i zielona (G) obiektu \texttt{Color}, tworząc gradient od zieleni dla małych wag do czerwieni dla dużych.
\end{itemize}

\section{Interakcja i Obsługa Zdarzeń}

\subsection{System Nawigacji}
Implementacja nawigacji wykorzystuje kombinację słuchaczy zdarzeń myszy:
\begin{itemize}
    \item \textbf{Zoom}: Obsługiwany przez \texttt{MouseWheelListener}. Skalowanie odbywa się względem aktualnej pozycji kursora myszy na ekranie.
    \item \textbf{Panning (Przesuwanie)}: Realizowane poprzez \texttt{MouseMotionListener}. Program zapamiętuje punkt startowy (\texttt{dragStart}) i aktualizuje przesunięcie (\texttt{px, py}) podczas przeciągania.
\end{itemize}

\subsection{Mechanizm Drag \& Drop}
Przesuwanie węzłów odbywa się wewnątrz \texttt{paintComponent}. Metoda \texttt{screenToWorld(Point p)} przelicza kliknięcie myszy na współrzędne układu grafu, uwzględniając aktualny zoom i przesunięcie, co pozwala na precyzyjne "chwycenie" węzła nawet przy dużym powiększeniu.

\section{Komunikacja z Silnikiem Zewnętrznym}
Wywołanie algorytmów zaimplementowanych w języku C (Fruchterman-Reingold / Tutte) odbywa się poprzez klasę \texttt{ProcessBuilder}. Mechanizm ten pozwala na uruchomienie skompilowanego pliku binarnego jako osobnego procesu systemowego.

Uruchomienie procesu w kodzie zrealizowano w następujący sposób:

\begin{lstlisting}[language=Java]
Process p = new ProcessBuilder(
    enginePath, 
    "-i", currentTopoFile, 
    "-o", "wynik.txt", 
    "-a", algo
).start();
\end{lstlisting}

Program przekazuje ścieżkę do pliku topologii (\texttt{-i}), nazwę pliku wyjściowego (\texttt{-o}) oraz wybrany przez użytkownika algorytm (\texttt{-a}). Następnie aplikacja oczekuje na zakończenie obliczeń metodą \texttt{p.waitFor()}, po czym automatycznie odświeża widok, wczytując wygenerowane współrzędne z pliku.
\section{Interfejs Użytkownika (UI Styling)}
W projekcie zrezygnowano z domyślnego wyglądu komponentów Swing na rzecz nowoczesnej estetyki:
\begin{itemize}
    \item \textbf{Custom Styling}: Metoda \texttt{styleButton} nadaje przyciskom kolory z palety pastelowej, usuwa obramowania i zmienia kursor na \texttt{HAND\_CURSOR}.
    \item \textbf{Layout}: Zastosowano \texttt{BoxLayout} dla panelu bocznego oraz \texttt{GridLayout} zamknięty w \texttt{CompoundBorder} dla widgetu właściwości węzła, co zapewnia czytelność i responsywność interfejsu.
\end{itemize}

\section{Specyfikacja plików danych}
\subsection{Format Topologii (.txt)}
Plik wejściowy czytany jest linia po linii. Każda linia musi zawierać co najmniej 4 kolumny rozdzielone białymi znakami:
\begin{verbatim}
[Nazwa_Krawędzi] [ID_Start] [ID_End] [Waga]
\end{verbatim}

\subsection{Format Współrzędnych (.txt)}
Zapis i odczyt realizowany jest zgodnie ze schematem:
\begin{verbatim}
[ID_Węzła] [X_double] [Y_double]
\end{verbatim}
Współrzędne zapisywane są z precyzją do 6 miejsc po przecinku przy użyciu \texttt{Locale.US} w celu zapewnienia kropki jako separatora dziesiętnego.

\end{document}